
\chapter{Results and discussions}

In this chapter, we examine how the kinematically cold and warm \HI\ components are distributed in nearby galaxies and how closely each component is linked to molecular gas and star formation.



\section[Spatial distribution of the kinematically cold and warm components]{Spatial distribution of the kinematically cold and warm components}
\label{sec:phase_spatial_distribution}

Figure~\ref{fig:individual_galaxies} shows the total, kinematically cold and kinematically warm \HI\ surface-density maps for the individual galaxies. Both components extend across much of the mapped atomic discs. Both components are present in clumps and ridges within the extended \HI\ emission.


The kinematically cold component contributes a median of about 45\% of the mapped \HI\ mass, while the warm component contributes the remaining 55\%. Within the optical disc ($r<r_{25}$), the warm component dominates the \HI\ budget, although the cold component remains a substantial part of the atomic gas.

The cold component has a more filamentary appearance than the warm component, with thin, elongated structures standing out against the more diffuse atomic emission. These elongated structures raise the question of whether the cold component is preferentially found in spiral arms. Such a preference could arise because compression in spiral arms can increase the gas density and promote cooling \citep{KimSpiral2008}. After tracing the spiral arms in nine spiral galaxies with visible arm structure in our sample following the methodology used by \citet{Foyle2010}, which used the 3.6 $\mu$m emission to trace stellar arms, we compared the fraction of the narrow component found in the arms and interarm regions in these galaxies. We found that the kinematically cold component is not consistently higher in the arm regions.
%To trace spiral arms in nine spiral galaxies with visible arm structure in our sample, we follow the methodology used by \citet{Foyle2010}, which used the 3.6 $\mu$m emission to trace stellar arms. We compare the fraction of the narrow component found in the arms and interarm regions in these galaxies. We found that the kinematically cold component is not consistently higher in the arm regions.

%maybe move the below to HI-H2 transition
Models of \HI\ thermal balance predict that cold, dense gas can coexist with warm, diffuse gas over a range of thermal pressures \citep{Wolfire2003}. At similar thermal pressure, the cold gas is denser because of its lower temperature. Simulations show that increasing the density of warm atomic gas can allow it to cool and condense into cold structures \citep{Saury2014}. Cold atomic gas can then supply material for molecular-cloud formation \citep{Krumholz2009Transition}.
%H$_2$ forms mainly on the surfaces of dust grains, and its formation is faster at higher gas densities \citep{GloverMacLow2007}. However, ultraviolet radiation can dissociate the molecules. Molecular gas can accumulate where dust and H$_2$ itself absorb enough of this radiation to protect the molecules. The transition from atomic to molecular gas therefore depends on the gas density, the shielding column and the dust abundance \citep{Krumholz2009Transition,Sternberg2014}.
In the simulations of \citet{GloverMacLow2007}, the initially atomic gas cools rapidly, and turbulent compression produces dense regions in which H$_2$ forms more rapidly. The connection to molecular-cloud formation therefore comes from the high densities that cold atomic gas can reach, together with sufficient shielding. This provides a reason to look for an association between the narrow \HI\ component and CO emission.

We found that both atomic components can be enhanced near CO-bright emission, but the cold fraction is not consistently higher there. The clearest enhancement occurs in NGC 5194, where the area-weighted mean cold fraction is about 45\% in the selected CO-bright regions, compared with 29\% in nearby CO-faint regions. In NGC 6946, both components increase near bright CO while the cold fraction remains about 30\%. The rest of the studied sample provides few suitable matched regions to enable a meaningful comparison between CO-bright and CO-faint emission. We therefore find no consistent evidence that narrow \HI\ is generally more concentrated near CO emission than broad \HI.

The absence of a consistent enhancement in the cold fraction does not rule out the fact that cold atomic gas is supplying material for molecular-cloud formation. In atomic-to-molecular transition models, atomic gas remains in shielding layers around molecular interiors, while additional gas can become molecular \citep{Krumholz2009Transition,Sternberg2014}. Once hydrogen becomes H$_2$, it no longer contributes to the \HI\ emission, so an increasing molecular reservoir need not be accompanied by an increasing narrow-\HI\ contribution. Furthermore, H$_2$ can survive where CO is dissociated \citep{Wolfire2010}. Our CO-faint comparison regions are therefore not necessarily free of H$_2$. These effects could weaken a spatial contrast between the CO-bright and CO-faint regions. Thus, our comparison does not establish whether the narrow \HI\ component preferentially traces regions where atomic gas becomes molecular. Locating the \HI-to-H$_2$ transition regions in our sample is difficult with our data, since a region may contain H$_2$ without detectable CO emission \citep{Wolfire2010}.


%



%\section[Which \HI\ component is more closely linked to star formation?]{Which \HI\ component is more closely linked to star formation?}
\section[Kennicutt--Schmidt relation Fitting methodology]{Kennicutt--Schmidt relation Fitting methodology}
\label{sec:ks_fitting}
\label{sec:sflaw}

%In this section, we test which of the kinematically cold and warm \HI\ components is more closely linked to star formation. We do so by fitting the Kennicutt--Schmidt relation for each gas component and comparing the slopes and correlation strengths, using total \HI\ and CO-traced H$_2$ as references. We investigate whether these relations differ between individual galaxies and between dwarf and spiral galaxies. We describe the fitting method in Section~\ref{sec:ks_fitting_method} and present the gas--SFR relations in Section~\ref{sec:ks_individual_results}.

%\subsection[Fitting methodology]{Fitting methodology}
\label{sec:ks_fitting_method}
To test which of the kinematically cold and warm \HI\ components is more closely linked to star formation, we fit the Kennicutt--Schmidt relation for each gas component and compare the slopes and correlation strengths, using total \HI\ and CO-traced H$_2$ as references. 
Following \citet{Fbigiel2008}, we compare the slopes of the Kennicutt--Schmidt relations for the different gas components. In logarithmic space, the Kennicutt--Schmidt relation in Equation~\ref{eq:SFLaw} becomes
\begin{equation}
\label{eq:ks_log_relation}
\log_{10}\Sigma_{\mathrm{SFR}} = a + N\log_{10}\Sigma_{\mathrm{gas}},
\end{equation}
where $a=\log_{10}A$ is the intercept and $N$ is the power-law index.

\citet{Fbigiel2008} used the ordinary least-squares (OLS) bisector to fit their Kennicutt--Schmidt relations. A single OLS regression minimises the vertical or horizontal residuals, depending on which variable is chosen as dependent. The bisector combines these two regressions and treats the axes symmetrically. Several atomic-gas distributions in our sample have considerable scatter or are nearly vertical. We therefore adopt orthogonal regression, which minimises the distances perpendicular to the fitted line and estimates the direction of the distribution directly.

Since neighbouring pixels within a beam contain correlated measurements of the same emission, we sample the maps at positions separated by at least one beam width, both horizontally and vertically. Before doing so, we first match the gas and SFR maps to the same spatial resolution and pixel grid. At each selected position, we use the gas and SFR values from the corresponding pixel.  We fit the Kennicutt--Schmidt relation to these beam-sampled measurements, shown as blue points in Figure~\ref{fig:individual_galaxies}. The grey shading shows all positive pixel values across the mapped field.

Each beam-sampled measurement in the Kennicutt--Schmidt relation has coordinates $(x_i,y_i)$, where $x_i=\log_{10}\Sigma_{\mathrm{gas},i}$ and $y_i=\log_{10}\Sigma_{\mathrm{SFR},i}$, for $i=1,\ldots,n$. Here, $n$ is the number of beam-sampled measurements. We use the orthogonal-regression criterion \citep{Pearson1901,Isobe1990}, with weights defined below. The fitted line minimises the weighted sum of squared perpendicular distances:
\begin{equation}
\label{eq:ks_weighted_distance}
\mathcal{Q}(a,N)
=\sum_{i=1}^{n} w_i\,d_{\perp,i}^{\,2}
=\sum_{i=1}^{n} w_i\,\frac{(y_i-a-Nx_i)^2}{1+N^2},
\end{equation}
where $N$ is the slope of the Kennicutt--Schmidt relation in Equation~\ref{eq:ks_log_relation}, $a$ is the intercept and $w_i$ is the weight assigned to measurement $i$. The distance $d_{\perp,i}$ is the shortest distance from $(x_i,y_i)$ to the fitted straight line $y=a+Nx$ on the plot of $\log_{10}\Sigma_{\mathrm{SFR}}$ against $\log_{10}\Sigma_{\mathrm{gas}}$.

To reduce the influence of sparsely distributed data points, we give greater weight to points in densely populated parts of the gas--SFR plot and less weight to points in sparse regions. We estimate the local point density using a two-dimensional Gaussian kernel \citep{Scott1992}. For each blue point $i$, we define its position in the logarithmic gas--SFR plot as
\begin{equation}
\mathbf{z}_i =
\begin{pmatrix}
x_i \\
y_i
\end{pmatrix}.
\end{equation}
We assign this point a weight proportional to the density of neighbouring blue points in the gas--SFR plot:
\begin{equation}
\label{eq:ks_density_weight}
w_i = \sum_{\substack{j=1\\j\ne i}}^{n}
\exp\!\left[-\frac{1}{2}
(\mathbf{z}_i-\mathbf{z}_j)^{\mathrm{T}}
\mathbf{H}^{-1}
(\mathbf{z}_i-\mathbf{z}_j)\right],
\end{equation}
where the sum includes all other blue points. The matrix $\mathbf{H}$ sets the width and shape of the Gaussian neighbourhood. We use Scott's bandwidth rule \citep{Scott1992}, giving $\mathbf{H}=n^{-1/3}\mathbf{S}$, where $n$ is the number of blue points and $\mathbf{S}$ is the covariance matrix of their logarithmic coordinates. Nearby points contribute strongly to the sum, while distant points contribute very little. A point with more nearby neighbours therefore receives a larger weight. These weights describe the concentration of measurements in the plot; they are not measurement-error weights.

We use the Pearson correlation coefficient $r$ to quantify the strength and direction of the correlation in each Kennicutt--Schmidt relation.

To estimate the uncertainties in the fitted parameters, we generate 2000 bootstrap samples by resampling the paired gas and SFR measurements with replacement. We fit the relation to each resampled data set. We use the 16th and 84th percentiles of the fitted values from these samples to define the lower and upper limits of the 68\% uncertainty interval. These intervals describe sampling uncertainty and exclude systematic uncertainties in the CO-to-H$_2$ conversion and SFR calibration. For nearly vertical relations where the slope is poorly constrained, we report the angle $\theta$ of the fitted line to the horizontal axis. The dashed lines in Figure~\ref{fig:individual_galaxies} show the fitted relations. Table~\ref{tab:individual_ks_summary} gives the fitted values and their uncertainties.


\begingroup
\raggedbottom
\setlength{\intextsep}{6pt}
%\section[Which \HI\ component is more closely linked to star formation?]{Which \HI\ component is more closely linked to star formation?}
\section[Kennicutt--Schmidt relations for the individual gas components]{Kennicutt--Schmidt relations for the individual gas components}
\label{sec:ks_individual_results}
To understand which \HI\ component is more closely linked to star formation, we first examine the Kennicutt--Schmidt relations in individual galaxies at their respective matched spatial resolutions. We then compare the dwarf and spiral samples by combining their measurements at a common physical resolution of 800 pc.\\\indent
The lower panels of Figure~\ref{fig:individual_galaxies} show the Kennicutt--Schmidt relations for total, narrow and broad \HI. Where available, H$_2$ is also shown for comparison. DDO 53 and DDO 154 are not shown because they provide only one and six beam-spaced measurements, respectively, within the adopted optical-disc boundary at 400 pc.

The atomic-gas relations vary considerably between galaxies (Table~\ref{tab:individual_ks_summary}). Ho II shows weak correlations with all three \HI\ tracers. In IC 2574 and NGC 0925, total \HI\ gives the strongest correlation with the SFR. Of the two separated components, narrow \HI\ correlates more strongly with the SFR than broad \HI. In NGC 2403, NGC 2976 and NGC 4214, total and broad \HI\ follow the SFR more closely than narrow \HI. NGC 7793 also shows a stronger correlation with total \HI\ than with either separated component. The relative strength of the narrow- and broad-\HI\ correlations with the SFR therefore varies between individual galaxies. This is consistent with \citet{Young2003}, who found no simple connection between the incidence of low-dispersion \HI\ and the SFR.

In the spiral galaxies, the narrow-component relations are weak in NGC 0628, NGC 3184 and NGC 6946. NGC 3351 shows a negative narrow-\HI\ correlation, while its broad component shows a weak positive correlation. Both atomic components show positive trends in NGC 2903 and NGC 4736, with the stronger relation for broad \HI. NGC 5194 also has a stronger broad than narrow relation, but total \HI\ gives the strongest atomic correlation. The strength and even the sign of the narrow-\HI\ correlation with the SFR differ between these spiral galaxies.

\citet{Fbigiel2008} also found that a small range of \HI\ surface density can accompany a much wider range of SFR surface density. They interpreted the saturation of the \HI\ surface density as evidence for a transition to molecular gas, with additional gas found mainly in the molecular phase. Under this interpretation, regions with similar \HI\ surface densities can have different molecular surface densities and SFRs.

These individual-galaxy comparisons use gas and SFR maps matched to the same spatial resolution within each galaxy. The physical size of the beam differs from one galaxy to another. The molecular relations are positive and are clearer than the atomic relations in most spirals. Their fitted slopes range from about 0.86 to 1.20. We therefore recover the approximately linear molecular relation already established by \citet{Fbigiel2008}, using density-weighted orthogonal regression instead of their OLS-bisector method. Here, the molecular relation provides a reference for assessing whether narrow or broad \HI\ is more closely associated with star formation. Among the lower-mass galaxies, NGC 2976 has a strong molecular correlation and a steeper slope, while NGC 2403 has a moderate, approximately linear relation. NGC 0925 and NGC 4214 have weaker molecular correlations and less well constrained slopes.

To examine the gas--SFR relations across the sample on the same physical scale, we smooth the maps to a common resolution of 800 pc and combine the beam-sampled measurements from the galaxies. We compare the dwarf and spiral samples separately, as well as the full sample. Figure~\ref{fig:pooled_dwarfs_ks} shows that the combined dwarf and spiral samples behave differently. In the combined dwarf sample, narrow \HI\ correlates more strongly with the SFR ($r=0.51$) than total \HI\ ($r=0.45$) or broad \HI\ ($r=0.23$). In spirals, narrow \HI\ gives the weakest relation ($r=0.11$), while broad and total \HI\ give $r=0.37$ and $r=0.34$. Combining the measurements from all galaxies gives $r=0.18$, 0.34 and 0.35 for narrow, broad and total \HI, respectively.

A possible explanation for the difference between dwarfs and spirals is that line width traces different gas conditions in the two populations. In dwarfs, narrow emission may preferentially trace cooler, less turbulent gas associated with conditions favourable for star formation. In spirals, cold gas near star-forming regions may have its emission broadened by turbulence. Such non-thermal motions can contribute to the line width of cold atomic gas \citep{Saury2014}. Energy input from young stars can help maintain \HI\ turbulence in star-forming discs \citep{Tamburro2009}. Some cold gas in these regions may therefore be assigned to the broad component, which could explain why the broad component shows stronger correlations with the SFR in spirals.


\begin{table}[p]
\centering
\caption{\textbf{Resolved gas--SFR fits for the individual galaxies and the combined sample.} $\theta$ denotes the angle between the fitted line and the horizontal axis and represents the steepness of the relation, with $N=\tan\theta$. For molecular gas, $N$ is the power-law index. The coefficient $r$ is the Pearson correlation coefficient.}
\label{tab:individual_ks_summary}
\par\vspace{8pt}
\begingroup

\footnotesize
\setlength{\tabcolsep}{3pt}
\renewcommand{\arraystretch}{1.12}
\textbf{(a) Atomic gas}\\[5pt]
\begin{tabular*}{\textwidth}{@{\extracolsep{\fill}}lrrrrrrr@{}}
\toprule
 & & \multicolumn{2}{c}{Total \HI} & \multicolumn{2}{c}{Narrow \HI} & \multicolumn{2}{c}{Broad \HI} \\
Galaxy & $n$ & $\theta$ [$^\circ$] & $r$ & $\theta$ [$^\circ$] & $r$ & $\theta$ [$^\circ$] & $r$ \\
\midrule
Ho II & 117 & $83.52^{+2.04}_{-2.06}$ & $0.26$ & $82.71^{+3.07}_{-2.40}$ & $0.18$ & $82.65^{+2.87}_{-2.90}$ & $0.23$ \\
IC 2574 & 455 & $79.94^{+1.09}_{-0.77}$ & $0.46$ & $72.12^{+2.08}_{-1.46}$ & $0.41$ & $82.87^{+1.28}_{-1.06}$ & $0.28$ \\
NGC 0925 & 380 & $81.87^{+0.79}_{-0.61}$ & $0.56$ & $80.96^{+1.45}_{-1.52}$ & $0.42$ & $83.23^{+1.33}_{-1.27}$ & $0.32$ \\
NGC 2403 & 382 & $79.46^{+1.11}_{-0.99}$ & $0.61$ & $82.90^{+1.11}_{-1.09}$ & $0.36$ & $74.00^{+1.19}_{-1.44}$ & $0.60$ \\
NGC 2976 & 97/96/97 & $69.94^{+0.88}_{-1.54}$ & $0.88$ & $67.44^{+0.67}_{-6.25}$ & $0.46$ & $70.31^{+1.30}_{-1.17}$ & $0.88$ \\
NGC 4214 & 88 & $79.50^{+2.06}_{-2.06}$ & $0.60$ & $81.61^{+4.12}_{-4.65}$ & $0.42$ & $76.84^{+2.15}_{-2.08}$ & $0.57$ \\
NGC 7793 & 352 & $78.26^{+0.79}_{-0.73}$ & $0.78$ & $74.29^{+1.55}_{-1.47}$ & $0.49$ & $79.80^{+0.97}_{-1.03}$ & $0.58$ \\
NGC 0628 & 428 & $86.99^{+0.84}_{-0.79}$ & $0.26$ & $-87.59^{+0.92}_{-1.19}$ & $-0.02$ & $81.44^{+1.40}_{-1.12}$ & $0.31$ \\
NGC 2903 & 181/177/175 & $74.83^{+1.28}_{-1.54}$ & $0.78$ & $70.02^{+2.68}_{-2.28}$ & $0.57$ & $75.92^{+1.54}_{-1.75}$ & $0.67$ \\
NGC 3184 & 611 & $84.89^{+0.50}_{-0.41}$ & $0.43$ & $-89.85^{+0.74}_{-0.81}$ & $0.04$ & $81.87^{+0.83}_{-0.80}$ & $0.44$ \\
NGC 3351 & 314 & $87.35^{+0.95}_{-0.51}$ & $0.01$ & $-85.74^{+1.33}_{-1.38}$ & $-0.27$ & $83.23^{+1.91}_{-1.52}$ & $0.18$ \\
NGC 4736 & 84 & $70.30^{+1.59}_{-0.93}$ & $0.93$ & $71.52^{+2.29}_{-1.21}$ & $0.84$ & $68.60^{+1.52}_{-1.05}$ & $0.93$ \\
NGC 5194 & 288 & $84.21^{+1.40}_{-1.28}$ & $0.42$ & $83.80^{+1.82}_{-1.39}$ & $0.26$ & $83.80^{+1.40}_{-1.37}$ & $0.38$ \\
NGC 6946 & 216 & $87.36^{+0.72}_{-0.92}$ & $0.36$ & $-89.83^{+2.09}_{-2.76}$ & $0.07$ & $86.83^{+1.15}_{-1.60}$ & $0.34$ \\
\midrule
Dwarfs & 769 & $83.21^{+0.53}_{-0.63}$ & $0.45$ & $80.26^{+0.74}_{-0.91}$ & $0.51$ & $85.23^{+0.80}_{-0.87}$ & $0.23$ \\
Spirals & 2119 & $84.77^{+0.47}_{-0.49}$ & $0.34$ & $89.83^{+0.51}_{-0.59}$ & $0.11$ & $78.49^{+0.70}_{-0.65}$ & $0.37$ \\
All galaxies & 2888 & $84.50^{+0.39}_{-0.39}$ & $0.35$ & $87.17^{+0.47}_{-0.49}$ & $0.18$ & $80.95^{+0.57}_{-0.52}$ & $0.34$ \\
\bottomrule
\end{tabular*}
\par\vspace{8pt}
\textbf{(b) Molecular gas}\\[5pt]
\begin{tabular*}{0.76\textwidth}{@{\extracolsep{\fill}}lrrr@{}}
\toprule
Galaxy & $n$ & $N$ & $r$ \\
\midrule
NGC 0925 & 257 & $2.12^{+0.86}_{-0.72}$ & $0.25$ \\
NGC 2403 & 259 & $0.92^{+0.05}_{-0.12}$ & $0.58$ \\
NGC 2976 & 86 & $1.55^{+0.12}_{-0.14}$ & $0.85$ \\
NGC 4214 & 66 & $2.32^{+0.60}_{-0.95}$ & $0.40$ \\
NGC 0628 & 387 & $1.05^{+0.04}_{-0.04}$ & $0.73$ \\
NGC 2903 & 146 & $1.02^{+0.04}_{-0.06}$ & $0.84$ \\
NGC 3184 & 522 & $1.20^{+0.05}_{-0.07}$ & $0.64$ \\
NGC 3351 & 264 & $1.15^{+0.09}_{-0.10}$ & $0.71$ \\
NGC 4736 & 79 & $1.20^{+0.06}_{-0.06}$ & $0.91$ \\
NGC 5194 & 276 & $0.90^{+0.11}_{-0.01}$ & $0.72$ \\
NGC 6946 & 212 & $0.86^{+0.05}_{-0.05}$ & $0.84$ \\
\midrule
Dwarfs & 507 & $1.15^{+0.04}_{-0.17}$ & $0.35$ \\
Spirals & 1918 & $1.07^{+0.02}_{-0.02}$ & $0.73$ \\
All galaxies & 2425 & $0.84^{+0.02}_{-0.02}$ & $0.60$ \\
\bottomrule
\end{tabular*}
\par\vspace{5pt}
\footnotesize\raggedright
$n$ is the number of fitted beam-spaced measurements.
\par
\endgroup
\end{table}

\FloatBarrier
\clearpage

\vspace*{\fill}
\begin{figure}[H]
\centering
\includegraphics[width=\linewidth]{images/HO_II-cnm_wnm.pdf}
\par\vspace{6pt}
\includegraphics[width=\linewidth]{images/ks_individual/HO_II-sfr_to_gas.pdf}
\caption[Surface-density maps and Kennicutt--Schmidt relations for individual galaxies]{\textbf{Ho II.} Top: total, kinematically cold, and kinematically warm \HI\ surface densities. Bottom: Kennicutt--Schmidt relations for total, narrow, and broad \HI, and H$_2$ where included. Grey shading shows all positive mapped pixels. Blue circles show gas and SFR measurements sampled at intervals of at least one beam width. Dashed lines show density-weighted orthogonal fits to these blue points (where available). For the relations between total, narrow or broad \HI\ and the SFR, $\theta$ is the angle between the fitted line and the horizontal axis. For the H$_2$--SFR relation, we report the slope $N$, or $\theta$ if the line is too close to vertical for a finite slope uncertainty to be determined. Intervals are central 68\% bootstrap intervals; $r$ is the Pearson coefficient and $n_{\mathrm{fit}}$ counts the H$_2$--SFR measurements used in the fit. Light grey lines mark constant depletion times.}
\label{fig:individual_galaxies}
\label{fig:hi_maps_HO_II}
\label{fig:ks_individual_HO_II}
\end{figure}
\vfill
\clearpage

\vspace*{\fill}
\begin{figure}[H]
\ContinuedFloat
\captionsetup{labelformat=continuedgalaxies,labelsep=none,list=no}
\centering
\textbf{IC 2574}\par\vspace{4pt}
\includegraphics[width=\linewidth]{images/IC2574-cnm_wnm.pdf}
\par\vspace{6pt}
\includegraphics[width=\linewidth]{images/ks_individual/IC2574-sfr_to_gas.pdf}
\caption[]{}
\label{fig:hi_maps_IC2574}
\label{fig:ks_individual_IC2574}
\end{figure}
\vfill
\clearpage

\vspace*{\fill}
\begin{figure}[H]
\ContinuedFloat
\captionsetup{labelformat=continuedgalaxies,labelsep=none,list=no}
\centering
\textbf{NGC 0925}\par\vspace{4pt}
\includegraphics[width=\linewidth]{images/NGC0925-cnm_wnm.pdf}
\par\vspace{6pt}
\includegraphics[width=\linewidth]{images/ks_individual/NGC0925-sfr_to_gas.pdf}
\caption[]{}
\label{fig:hi_maps_NGC0925}
\label{fig:ks_individual_NGC0925}
\end{figure}
\vfill
\clearpage

\vspace*{\fill}
\begin{figure}[H]
\ContinuedFloat
\captionsetup{labelformat=continuedgalaxies,labelsep=none,list=no}
\centering
\textbf{NGC 2403}\par\vspace{4pt}
\includegraphics[width=\linewidth]{images/NGC2403-cnm_wnm.pdf}
\par\vspace{6pt}
\includegraphics[width=\linewidth]{images/ks_individual/NGC2403-sfr_to_gas.pdf}
\caption[]{}
\label{fig:hi_maps_NGC2403}
\label{fig:ks_individual_NGC2403}
\end{figure}
\vfill
\clearpage

\vspace*{\fill}
\begin{figure}[H]
\ContinuedFloat
\captionsetup{labelformat=continuedgalaxies,labelsep=none,list=no}
\centering
\textbf{NGC 2976}\par\vspace{4pt}
\includegraphics[width=\linewidth]{images/NGC2976-cnm_wnm.pdf}
\par\vspace{6pt}
\includegraphics[width=\linewidth]{images/ks_individual/NGC2976-sfr_to_gas.pdf}
\caption[]{}
\label{fig:hi_maps_NGC2976}
\label{fig:ks_individual_NGC2976}
\end{figure}
\vfill
\clearpage

\vspace*{\fill}
\begin{figure}[H]
\ContinuedFloat
\captionsetup{labelformat=continuedgalaxies,labelsep=none,list=no}
\centering
\textbf{NGC 4214}\par\vspace{4pt}
\includegraphics[width=\linewidth]{images/NGC4214-cnm_wnm.pdf}
\par\vspace{6pt}
\includegraphics[width=\linewidth]{images/ks_individual/NGC4214-sfr_to_gas.pdf}
\caption[]{}
\label{fig:hi_maps_NGC4214}
\label{fig:ks_individual_NGC4214}
\end{figure}
\vfill
\clearpage

\vspace*{\fill}
\begin{figure}[H]
\ContinuedFloat
\captionsetup{labelformat=continuedgalaxies,labelsep=none,list=no}
\centering
\textbf{NGC 7793}\par\vspace{4pt}
\includegraphics[width=\linewidth]{images/NGC7793-cnm_wnm.pdf}
\par\vspace{6pt}
\includegraphics[width=\linewidth]{images/ks_individual/NGC7793-sfr_to_gas.pdf}
\caption[]{}
\label{fig:hi_maps_NGC7793}
\label{fig:ks_individual_NGC7793}
\end{figure}
\vfill
\clearpage

\vspace*{\fill}
\begin{figure}[H]
\ContinuedFloat
\captionsetup{labelformat=continuedgalaxies,labelsep=none,list=no}
\centering
\textbf{NGC 0628}\par\vspace{4pt}
\includegraphics[width=\linewidth]{images/NGC0628-cnm_wnm.pdf}
\par\vspace{6pt}
\includegraphics[width=\linewidth]{images/ks_individual/NGC0628-sfr_to_gas.pdf}
\caption[]{}
\label{fig:hi_maps_NGC0628}
\label{fig:ks_individual_NGC0628}
\end{figure}
\vfill
\clearpage

\vspace*{\fill}
\begin{figure}[H]
\ContinuedFloat
\captionsetup{labelformat=continuedgalaxies,labelsep=none,list=no}
\centering
\textbf{NGC 2903}\par\vspace{4pt}
\includegraphics[width=\linewidth]{images/NGC2903-cnm_wnm.pdf}
\par\vspace{6pt}
\includegraphics[width=\linewidth]{images/ks_individual/NGC2903-sfr_to_gas.pdf}
\caption[]{}
\label{fig:hi_maps_NGC2903}
\label{fig:ks_individual_NGC2903}
\end{figure}
\vfill
\clearpage

\vspace*{\fill}
\begin{figure}[H]
\ContinuedFloat
\captionsetup{labelformat=continuedgalaxies,labelsep=none,list=no}
\centering
\textbf{NGC 3184}\par\vspace{4pt}
\includegraphics[width=\linewidth]{images/NGC3184-cnm_wnm.pdf}
\par\vspace{6pt}
\includegraphics[width=\linewidth]{images/ks_individual/NGC3184-sfr_to_gas.pdf}
\caption[]{}
\label{fig:hi_maps_NGC3184}
\label{fig:ks_individual_NGC3184}
\end{figure}
\vfill
\clearpage

\vspace*{\fill}
\begin{figure}[H]
\ContinuedFloat
\captionsetup{labelformat=continuedgalaxies,labelsep=none,list=no}
\centering
\textbf{NGC 3351}\par\vspace{4pt}
\includegraphics[width=\linewidth]{images/NGC3351-cnm_wnm.pdf}
\par\vspace{6pt}
\includegraphics[width=\linewidth]{images/ks_individual/NGC3351-sfr_to_gas.pdf}
\caption[]{}
\label{fig:hi_maps_NGC3351}
\label{fig:ks_individual_NGC3351}
\end{figure}
\vfill
\clearpage

\vspace*{\fill}
\begin{figure}[H]
\ContinuedFloat
\captionsetup{labelformat=continuedgalaxies,labelsep=none,list=no}
\centering
\textbf{NGC 4736}\par\vspace{4pt}
\includegraphics[width=\linewidth]{images/NGC4736-cnm_wnm.pdf}
\par\vspace{6pt}
\includegraphics[width=\linewidth]{images/ks_individual/NGC4736-sfr_to_gas.pdf}
\caption[]{}
\label{fig:hi_maps_NGC4736}
\label{fig:ks_individual_NGC4736}
\end{figure}
\vfill
\clearpage

\vspace*{\fill}
\begin{figure}[H]
\ContinuedFloat
\captionsetup{labelformat=continuedgalaxies,labelsep=none,list=no}
\centering
\textbf{NGC 5194}\par\vspace{4pt}
\includegraphics[width=\linewidth]{images/NGC5194-cnm_wnm.pdf}
\par\vspace{6pt}
\includegraphics[width=\linewidth]{images/ks_individual/NGC5194-sfr_to_gas.pdf}
\caption[]{}
\label{fig:hi_maps_NGC5194}
\label{fig:ks_individual_NGC5194}
\end{figure}
\vfill
\clearpage

\vspace*{\fill}
\begin{figure}[H]
\ContinuedFloat
\captionsetup{labelformat=continuedgalaxies,labelsep=none,list=no}
\centering
\textbf{NGC 6946}\par\vspace{4pt}
\includegraphics[width=\linewidth]{images/NGC6946-cnm_wnm.pdf}
\par\vspace{6pt}
\includegraphics[width=\linewidth]{images/ks_individual/NGC6946-sfr_to_gas.pdf}
\caption[]{}
\label{fig:hi_maps_NGC6946}
\label{fig:ks_individual_NGC6946}
\end{figure}
\vfill
\clearpage

\vspace*{\fill}
\begin{figure}[H]
\centering
\includegraphics[width=\linewidth]{images/ks_individual/pooled_groups_ks.pdf}
\caption[Combined Kennicutt--Schmidt relations for dwarfs, spirals and all galaxies]{Combined Kennicutt--Schmidt relations for dwarfs (top), spirals (middle) and all galaxies (bottom), at a common physical resolution of 800 pc. \protect The columns show total, kinematically cold and kinematically warm \HI, and H$_2$, from left to right. Symbols and density-weighted fits follow Figure \ref{fig:individual_galaxies}. Fits to the \HI--SFR relations show $\theta$, and fits to the H$_2$--SFR relation show $N$.}
\label{fig:pooled_dwarfs_ks}
\end{figure}
\vfill
\clearpage
\FloatBarrier

\endgroup



\begingroup
\raggedbottom
\setlength{\parskip}{6pt}
\section[Star formation efficiency]{Star formation efficiency}
\label{sec:sfe}

The Kennicutt--Schmidt relations show how the gas and SFR surface densities vary together. Following \citet{Leroy2008}, we ask how the apparent SFE differs between the narrow and broad \HI\ reservoirs, and how it changes across the disc.% We define the star formation efficiency (SFE) as the SFR divided by the gas mass.

We calculate the SFE from the new narrow and broad maps, using the same matched spatial resolution within each galaxy as in the individual Kennicutt--Schmidt comparisons. For a region containing $N_{\mathrm{reg}}$ valid pixels, we use
\begin{equation}
\label{eq:regional_sfe}
\mathrm{SFE}_p = 10^{-6}\,
\frac{\displaystyle\sum_{i=1}^{N_{\mathrm{reg}}}\Sigma_{\mathrm{SFR},i} A_i}
{\displaystyle\sum_{i=1}^{N_{\mathrm{reg}}}\Sigma_{p,i} A_i},
\end{equation}
where $A_i$ is the pixel area and $p$ denotes the gas tracer. The factor $10^{-6}$ converts the SFR surface density from per kpc$^2$ to per pc$^2$, giving SFE in yr$^{-1}$. To compare the tracers over the same area, we calculate their SFEs from the same pixel values in the gas and SFR maps.

\begin{table}[tbp]
\centering
\caption{\textbf{Median aperture $\operatorname{\mathbf{log}}_{\mathbf{10}}\mathbf{SFE}$,} in yr$^{-1}$, calculated separately for each galaxy before combining the sample. Intervals are the central 68\% of 10\,000 whole-galaxy bootstrap resamples; calibration and decomposition uncertainties are not included. Atomic rows use nine dwarfs and seven spirals. Molecular rows use the same CO-covered pixels within four dwarfs and seven spirals.}
\label{tab:sfe}
\begingroup

\renewcommand{\arraystretch}{1.3}
\begin{tabular}{lrrr}
\toprule
Gas denominator & Dwarfs & Spirals & All galaxies \\
\midrule
Total \HI & $-9.69^{+0.21}_{-0.10}$ & $-9.16^{+0.17}_{-0.13}$ & $-9.39^{+0.14}_{-0.29}$ \\
Narrow \HI & $-9.29^{+0.21}_{-0.15}$ & $-8.66^{+0.05}_{-0.33}$ & $-9.08^{+0.19}_{-0.28}$ \\
Broad \HI & $-9.47^{+0.20}_{-0.06}$ & $-8.99^{+0.24}_{-0.00}$ & $-9.13^{+0.14}_{-0.27}$ \\
H$_2$ & $-8.45^{+0.19}_{-0.09}$ & $-9.10^{+0.13}_{-0.13}$ & $-8.96^{+0.39}_{-0.15}$ \\
H$_2$ + narrow \HI & $-9.02^{+0.13}_{-0.23}$ & $-9.25^{+0.02}_{-0.10}$ & $-9.23^{+0.17}_{-0.12}$ \\
H$_2$ + broad \HI & $-8.99^{+0.04}_{-0.27}$ & $-9.35^{+0.09}_{-0.03}$ & $-9.26^{+0.25}_{-0.13}$ \\
H$_2$ + total \HI & $-9.22^{+0.06}_{-0.28}$ & $-9.43^{+0.01}_{-0.02}$ & $-9.42^{+0.20}_{-0.03}$ \\
\bottomrule
\end{tabular}
\endgroup
\end{table}

Figure~\ref{fig:adopted_sfe} shows the SFE as a function of radius for the different gas components. Each point represents the median SFE of the galaxies contributing to a radial bin of width $0.1r_{25}$. In dwarfs, the narrow \HI\ reservoir appears more efficient at forming stars than the broad \HI\ reservoir throughout the disc. In spirals, narrow \HI\ has the higher SFE in the inner disc ($r<0.5r_{25}$), while the efficiencies of the narrow and broad components become more similar farther out. In both samples, the apparent SFE decreases towards the outer disc for total \HI\ and for either atomic component, consistent with the radial decline found in \HI-dominated regions by \citet{Leroy2008}.

The molecular SFE is comparatively flat with radius in spirals, while adding either atomic component introduces a radial decline. Adding \HI\ also lowers the SFE values. The H$_2$+narrow and H$_2$+broad profiles are similar.

In the innermost regions of spirals, the H$_2$+narrow and H$_2$+broad SFEs are close to the molecular-only SFE because the gas in these CO-covered regions is predominantly molecular. Adding either atomic component therefore makes little difference. Farther out, the atomic contribution increases, so adding \HI\ lowers the SFE more and the combined profiles diverge from the molecular-only profile, consistent with the radial decrease in molecular fraction found by \citet{Leroy2008}. In the CO-covered dwarfs, the H$_2$+narrow and H$_2$+broad SFEs are lower than the molecular-only SFE even in the central regions. Unlike spiral centres, these regions already contain a substantial atomic-gas contribution. As this contribution increases outwards, the two combined profiles fall farther below the molecular-only profile.

Table~\ref{tab:sfe} summarises the median SFEs of the dwarf and spiral samples.

\begin{figure}[H]
\centering
\includegraphics[width=\linewidth]{images/ks_individual/adopted_sfe_profiles.pdf}
\caption[Star formation efficiency with the new kinematic component maps]{Radial SFE profiles for dwarfs (top) and spirals (bottom). Left: the observed SFR divided by total, narrow or broad \HI. Right: H$_2$ and its combinations with either atomic component or total \HI, over identical CO-covered positions. Each point is the median of the galaxy measurements in a $0.1r_{25}$ bin. Shading shows the 16th--84th percentile spread between galaxies for total \HI\ (left) and H$_2$ (right), rather than uncertainty on the median.}
\label{fig:adopted_sfe}
\end{figure}
\FloatBarrier
\clearpage
\endgroup

